\documentclass[11pt,a4paper]{article}
\usepackage[T1]{fontenc}
\usepackage[utf8]{inputenc}
\usepackage[margin=2.3cm]{geometry}
\usepackage{xcolor}
\usepackage{enumitem}
\usepackage{hyperref}
\setlength{\parskip}{0.5em}
\setlength{\parindent}{0pt}
\newcommand{\referee}[1]{\par\medskip\noindent{\color{black!60}\itshape #1}\par}
\newcommand{\reply}{\noindent\textbf{Reply.}~}

\begin{document}

\begin{center}
{\large\bfseries Response to the Referees}\\[3pt]
Manuscript ZV10235: ``Accessing the symmetric component of quadrupole misalignment through
counter-rotating closed orbits in EDM storage rings''\\
I.~Abedrabbo and S.~Hac{\i}\"omero\u{g}lu
\end{center}

We thank both referees for their careful reading and constructive comments. The comments led us to test
the method against the error sources raised by Referee~A and to rewrite the presentation along the lines
requested by Referee~B. The main changes are:
\begin{enumerate}[nosep]
\item A new section ``Domain of validity'' (Sec.~V) quantifies the scheme against monitor motion,
electrostatic deflector field errors and rolls, corrector type and correlated misalignment.
\item A new section ``Operating scenario'' (Sec.~VI) follows the machine for one week with ground motion,
BPM drift and periodic BBA.
\item The two-beam scheme is now stated as an explicit algorithm (Sec.~IV.B), including the treatment of
BPM offsets and the choice of the regularization.
\item The symmetric and antisymmetric components are defined in the introduction (Sec.~I.C), and the claim
of a 100-fold orbit suppression of the symmetric component has been corrected: the measured suppression is
four- to sevenfold at low harmonics, explained by the different beta functions at the focusing and
defocusing quadrupoles (Sec.~III, Fig.~2).
\item The false EDM is now evaluated at the frozen-spin energy (Sec.~II, Appendix~A). With a launch at
fixed momentum, the result depended on the launch point in a misaligned lattice; this affected the
uncorrected and single-beam results by factors of 1.1--35, but not the two-beam results, which were
and remain at the numerical floor. The false-EDM values in the paper were recomputed accordingly.
\textcolor{red}{[Authors: BBA convergence and channel figure are being re-run with the new evaluation;
update this item and the figure numbers once they are included.]}
\item To keep the length under control, the main text was streamlined, the BBA section shortened, and
technical material moved to appendices.
\end{enumerate}

%======================================================================
\section*{Referee A}
%======================================================================

\referee{First, the paper assumes that the orbit perturbation is mainly driven by misalignment of quadrupoles
with alignment accuracy of about 50--100~$\mu$m. [\dots] there are other sources of orbit perturbation.
They are errors of bending power of electrostatic bends and their roles. Optimistically one can put the
relative accuracy of about $10^{-4}$ and the roll of 0.1~mrad. Then for the suggested optics, the orbit
perturbation from just one of 48 bends is about 400~$\mu$m.}

\reply We agree that this is essential, and we have added Sec.~V.B. We confirm the estimate: a relative field
error of $10^{-4}$ (roll of 0.1~mrad) in one deflector produces a closed-orbit distortion of 418~$\mu$m
(525~$\mu$m) peak. Such static errors are removed in the initial calibration; the relevant question for our
scheme is the \emph{drift} of the deflectors between calibrations, because the electric force does not
reverse with the beam direction and the deflector-induced orbit therefore appears in the same (sum) channel
as the symmetric quadrupole misalignment. We find:
\begin{itemize}[nosep]
\item A fit containing only quadrupole columns mistakes a deflector drift for a symmetric misalignment
(about 1~$\mu$m per $10^{-7}$ field drift or 0.1~$\mu$rad roll drift). Adding the deflector columns to the
fit removes this leakage, at the price of a larger noise amplification [Fig.~6(a)].
\item A deflector roll produces a first-order false EDM through the magnetic kicks that close the orbit. We
show that it depends only on the mean roll (zero-mean roll patterns give no first-order signal) and that it lives
in the even CW+CCW channel while the EDM lives in the odd channel. The common vertical electric trim already
foreseen by Omarov \emph{et al.} [PRD \textbf{105}, 032001 (2022), Sec.~III~B, Table~IV], which keeps the same-sign precession below $10^{-6}$~rad/s,
therefore removes it; we verify that this trim leaves the EDM untouched (Table~III). A CW/CCW energy asymmetry
$\Delta\delta$ leaks the residual even signal into the odd channel as $C\approx|S_+|\cdot2.7\,\Delta\delta$; with
that trim level the leak is negligible ($C\sim3\times10^{-4}$ for $\Delta\delta=10^{-7}$).
\item Zero-mean rolls still couple to the closed orbit at second order, in the odd channel; individual
deflector trims remove this term (matched test: $C=1.7$ with the common trim alone, $0.010$ with common and
individual trims).
\item Field errors and rolls together produce a geometric phase of their own. It is removed by trimming the
deflectors at the source (field by the deflector voltage, roll by an up--down electrode voltage difference),
using the deflector part of the fit. The tolerable deflector drift between trims is set by the accuracy of
the response model: with an exact model the false EDM stays at the numerical floor up to $10^{-5}$; with a
0.2\% gradient error the target is reached at $3$--$5\times10^{-6}$ [Fig.~6(b)].
\end{itemize}

\referee{The second problem is that the authors assume that accuracy of quad positioning is independent on
their location in the ring. [\dots] A more accurate model should account that an error in relative
positions of two quads is proportional to the square root of distance between quads. [\dots] The problem
is complicated by suppression of beam orbit deviation from BPM centers for long wave perturbations, but it
is not obvious that spin precession out of plane has similar suppression.}

\reply We have added Sec.~V.D with a random-walk (ATL-type) model in which the relative displacement grows
as the square root of the distance. Because neighboring quadrupoles move together, the symmetric
component dominates (84\% of the variance after removal of the rigid translation and tilt, against about
50\% for uncorrelated errors). The two-beam correction is linear and its residual does not depend on the
pattern statistics. The time-dependent version of this model is used in the one-week scenario of
Sec.~VI.

On the second point---whether the spin shares the suppression of long-wavelength perturbations---our answer
is yes, for the components that the beam actually follows. A uniform translation or tilt of the lattice
moves the closed orbit with it, so the beam still passes through the quadrupole centers and the spin is not
affected. More generally, the spin is most sensitive near the betatron tune ($m=2$--3) and weakly sensitive
at $m=0$. We also found that the more important consequence of long-wavelength motion is that the BPMs
move together with the quadrupoles. Section~V.A treats this explicitly: an estimator that assumes fixed
BPMs then recovers only about a third of the symmetric drift, but the remaining part lies at harmonics to
which the spin is weakly sensitive and the false EDM stays below target with a margin of three to five.
Monitors rigidly attached to the quadrupoles, with an estimator that accounts for the mounting, recover it
fully and are robust against BPM--quadrupole drifts up to at least 1~$\mu$m.

\referee{It also should be spelled out the difference between usage of electrostatic and magnetic
correctors.}

\reply Added as Sec.~V.C. Only dipole correctors mounted in the quadrupoles (or quadrupole movers) cancel
the misalignment where it occurs and keep the false EDM at the floor. Magnetic correctors at the deflectors
act only on the difference channel and cannot correct a symmetric drift. Electrostatic correctors act on
both beams alike and can steer them through the quadrupole centers, but the electric kicks they need
generate first- and second-order deflector effects ($C=1.3$--$7.6$ times the target, plus same-sign
signals that rely on CW/CCW symmetry).

\referee{Overall impression is that the paper is too long for the reported context.}

\reply We have rewritten the paper concisely: although new sections were added to address the points
above and those of Referee~B, the revised manuscript is shorter than the original (8 instead of 9 pages).
Technical material was moved to appendices.

\referee{1. Excessive use of acronyms greatly complicates reading of the text.}

\reply We removed most acronyms. The remaining ones (EDM, BPM, BBA, SVD, CW/CCW) are defined at first use.
\textcolor{red}{[Authors: final pass on acronyms.]}

\referee{2. Section 1.A, line 1: ``A permanent electric dipole moment violates'' $\to$ ``A permanent electric
dipole moment of an elementary particle violates''.}

\reply Changed as suggested.

\referee{3. Section 1C, line 6: ``Subtracting the'' $\to$ ``Subtraction of the''.}

\reply The sentence has been rewritten.

\referee{4. Section 1D: An explanation of symmetric and antisymmetric quad displacement comes late in the
text.}

\reply The definitions now appear in Sec.~I.C, before they are first used.

\referee{5. Section II.A, line 3: ``the lattice as described in Sec.~IB'' $\to$ ``the lattice described in
Sec.~IB''.}

\reply Changed.

\referee{6. Line before Eq.~(9): $\Delta R$ does not look as a matrix.}

\referee{7. Equation after (9) has problem with dimension.}

\reply (6, 7) The orbit-difference discussion was moved to Appendix~B and rewritten: $\Delta R$ is the
difference of the response matrices for two gradient settings, $\Delta R\approx\epsilon_g\,\partial
R/\partial g$, and the relation now reads $\|\Delta R^{-1}\|\approx\|R^{-1}\|/\epsilon_g$, with $\epsilon_g$
the dimensionless relative gradient change.

\referee{8. Section V.A, Line 3: In the phrase ``beams pass through the quadrupoles in the opposite order,''
state that these are two nearby quads.}

\reply Clarified: ``the reversed order of the two quadrupoles of a cell'' (Sec.~IV.B).

%======================================================================
\section*{Referee B}
%======================================================================

\referee{A hypothetical ring is described in Section 1.B. But no corresponding lattice layout or optics plot is
provided.}

\reply Added as Fig.~1 (cell layout and beta functions), together with the beta functions at the
quadrupoles in Table~I.

\referee{The manuscript introduces the symmetric and antisymmetric quadrupole modes near the beginning of the
article, but their mathematical definitions do not appear until Section III.A, in Equation (4). In
addition, Equation (5) uses $n$ as an index, whereas $k$ is not defined. [\dots]}

\reply The definitions now appear in Sec.~I.C, Eq.~(3), with the indices defined (quadrupole $n=0,\dots,47$,
cell $k=0,\dots,23$, QF $n=2k$, QD $n=2k+1$), the decomposition $v[n]=v^s[n]+v^a[n]$, and the statement
that the two projections are orthogonal and span 24-dimensional subspaces.

\referee{The authors state that the low-order singular-value-decomposition (SVD) modes correspond to the
symmetric mode, but they provide no explanation or supporting evidence for this correspondence.}

\reply Figure~2(b) now shows the singular values of the tracked response matrix colored by the symmetric
content of the corresponding right singular vectors: the directions with the smallest singular values are
predominantly symmetric. Section~III.A explains why: the symmetric pattern kicks the beam in opposite
directions at the two quadrupoles of a cell, and these kicks cancel up to the difference of
$\sqrt{\beta}$ between the two quadrupoles; the net kick is about 0.15 of that of an antisymmetric
pattern, in agreement with the measured orbit gains [Fig.~2(a)]. In revising this point we found that the
statement of a 100-fold suppression in the original manuscript, based on a smooth-focusing estimate, was
incorrect; the suppression is four- to sevenfold at low harmonics.

\referee{The central purpose of the paper is to present a method for keeping false EDM signals arising from
quadrupole misalignments below the sub-target threshold by monitoring the sum of the BPM readings from the
two beams. However, the implementation of this method is described only verbally in a single paragraph in
Section V.B. [\dots] a PRAB article requires a more rigorous presentation demonstrating how the method works
and how its performance is evaluated.}

\reply We agree. Section~IV.C now states the scheme as an explicit five-step algorithm: reference readings
after BBA, differences of both beams relative to the reference, the two-beam least-squares estimate
[Eq.~(4)] with the regularization chosen by the discrepancy principle (no knowledge of the true
misalignment is needed), the correction by quadrupole-mounted dipole correctors, and the BBA trigger. The
performance is evaluated in three steps: (i) static correction of random misalignments (Fig.~4);
(ii) the validity domain against monitor motion, deflector errors, corrector type and correlated
misalignment (Sec.~V, Figs.~5 and 6); (iii) a one-week time-dependent scenario with ground motion, BPM
drift, BPM noise and periodic BBA, in which two-beam monitoring keeps the false EDM at the level set by the
BBA residual (median 0.05--0.09 of the target, maximum 0.18), whereas periodic BBA alone rises to about
45 times the target between calibrations (Sec.~VI, Fig.~7). The detection limit of the monitor is given in
the same section.

\referee{Recommendation: the manuscript undergo major revision before it can be considered further for
publication.}

\reply We hope that the revised manuscript addresses the referees' concerns.

\bigskip
\textbf{Additional change not requested by the referees.} While revising, we found that our false-EDM
estimator launched the particle on the closed orbit at fixed momentum. With hard-edge deflectors the
momentum varies along a distorted closed orbit, so the result depended on the launch point. Since the
experiment keeps the horizontal spin frozen by feedback, we now evaluate the false EDM at the frozen-spin
energy (Sec.~II, Appendix~A), where the result is independent of the launch point to 0.2\%. The two-beam
results of the original manuscript are unaffected (both at the numerical floor); the uncorrected and
single-beam values decreased by factors of 1.1--35. All values in the revised manuscript use the new
evaluation.

\end{document}
